\begin{abstract}
The intricacies and formal structures inherent in mathematical reasoning make it a formidable task for language models. Here, we present DeepSeekMath~7B, an extension of DeepSeek-Coder-Base-v1.5 7B, further pre-trained on a vast corpus of 120B math-focused tokens extracted from Common Crawl, supplemented with code and natural language materials. Notably, DeepSeekMath~7B attains a notable 51.7\% on the competition-grade MATH benchmark, all without dependence on auxiliary toolkits or ensemble strategies, attaining results on par with Gemini-Ultra and GPT-4. Evaluated with self-consistency across 64 samples, DeepSeekMath~7B reaches a score of 60.9\% on MATH. The model's enhanced mathematical reasoning is primarily ascribed to two elements: firstly, the effective exploitation of large-scale, publicly sourced web data enabled by a carefully crafted data curation system; secondly, the application of Group Relative Policy Optimization (GRPO)—an adaptation of Proximal Policy Optimization (PPO)—which not only augments reasoning capabilities but also reduces PPO's memory footprint.
\end{abstract}

\section{Introduction}

Large language models (LLMs) have dramatically altered how artificial intelligence approaches mathematical reasoning, catalyzing substantial progress in quantitative reasoning tasks~\citep{MATH} as well as in benchmarks assessing geometric reasoning~\citep{trinh2024solving}. Furthermore, these systems have shown efficacy in supporting human problem-solving for complex mathematical challenges~\citep{tao}. Despite these achievements, state-of-the-art solutions such as GPT-4~\citep{gpt4} and Gemini-Ultra~\citep{gemini} remain proprietary, while current open-source alternatives lag significantly in their performance capabilities.

In this work, we present DeepSeekMath, a specialized language model focused on mathematical tasks that markedly surpasses other open-source models in mathematical reasoning and comes close to matching the proficiency of GPT-4 on standard academic assessments. To realize this, we construct the DeepSeekMath Corpus—a substantial, meticulously curated pre-training dataset containing 120B mathematical tokens. The corpus is assembled from Common Crawl (CC) utilizing a classifier based on fastText \citep{joulin2016fasttext}. During the first phase, the classifier is trained with positive examples drawn from OpenWebMath \citep{openwebmath} and a wide variety of other web content as negative examples. The classifier is then leveraged to identify additional positive samples within CC, which are subsequently validated and refined with the help of human annotators. Following this, the classifier is retrained with the augmented and improved dataset to achieve better accuracy. Experimental results demonstrate that the corpus maintains a high level of quality: our base model, DeepSeekMath-Base 7B, scores 64.2\% on GSM8K \citep{gsm8k} and 36.2\% on the MATH competition benchmark \citep{MATH}, outperforming models such as Minerva 540B \citep{minerva}. Moreover, the DeepSeekMath Corpus supports multiple languages, which is reflected in enhanced performance on Chinese mathematical evaluations \citep{wei2023cmath,agieval}. We contend that our methodological approach to mathematical data curation provides a foundation for further progress in the field, with ample opportunities for future enhancements.

DeepSeekMath-Base is built upon DeepSeek-Coder-Base-v1.5 7B \citep{deepseek-coder}, as we have found that initializing with a code-centric pre-trained model yields superior results compared to using a general-purpose large language model. Additionally, our findings suggest that math-focused pre-training not only augments mathematical performance but also bolsters abilities on general benchmarks such as MMLU \citep{mmlu} and BBH \citep{bbh}, demonstrating improvements in broader reasoning skills.

Following the pre-training phase, we implement mathematical instruction tuning on DeepSeekMath-Base utilizing datasets that encourage chain-of-thought reasoning \citep{cot}, program-of-thought methods \citep{pot,pal}, and tool-augmented reasoning \citep{tora}. The resultant model, DeepSeekMath-Instruct 7B, surpasses all other 7B models and reaches performance levels close to leading 70B open-source instruction-tuned systems.

In addition, we propose Group Relative Policy Optimization (GRPO), a reinforcement learning (RL) technique inspired by Proximal Policy Optimization (PPO) \citep{schulman2017proximal}. Distinctively, GRPO omits the traditional critic model and instead estimates baselines through group-level scoring, substantially lowering computational requirements. With access to only a segment of English instruction tuning datasets, GRPO delivers marked advancements over the robust DeepSeekMath-Instruct, both on in-domain tasks (GSM8K: 82.9\% $\rightarrow$ 88.2\%, MATH: 46.8\% $\rightarrow$ 51.7\%) and out-of-domain tasks such as CMATH (84.6\% $\rightarrow$ 88.8\%) during the RL stage. We further establish a unified conceptual framework that encompasses various methods, including Rejection Sampling Fine-Tuning (RFT) \citep{yuan2023scaling}, Direct Preference Optimization (DPO) \citep{dpo}, PPO, and GRPO. This framework reveals that these techniques can be understood as forms of direct or simplified reinforcement learning. We conduct comprehensive evaluations—contrasting online and offline training, outcome versus process supervision, and single-turn with iterative RL—to analyze key factors underpinning this approach. Finally, we elucidate how our RL methodology enhances the effectiveness of instruction-tuned models and highlight promising future directions for more efficient RL grounded in this unified view.

\subsection{Contributions}
This work introduces scalable approaches to mathematical pre-training, alongside an in-depth investigation into reinforcement learning techniques.

\textbf{Math Pre-Training at Scale}

\begin{itemize}[topsep=0pt]
    \item We demonstrate that the Common Crawl dataset, which is freely available to the public, serves as a rich resource for mathematical data. Through the deployment of a rigorous data selection pipeline, we assemble the DeepSeekMath Corpus—a curated collection containing 120B tokens extracted from web sources identified as mathematically relevant. This dataset surpasses the math-related data utilized in Minerva \citep{minerva} by nearly sevenfold, and is approximately nine times larger than the newly introduced OpenWebMath \citep{openwebmath}.

    \item DeepSeekMath-Base 7B, our foundational model trained on this dataset, exhibits performance on par with the much larger Minerva 540B \citep{minerva}, suggesting that parameter count is not the sole determinant of mathematical reasoning proficiency. Robust mathematical capabilities can emerge from more compact models given carefully curated training data.

    \item The results from our series of mathematical training experiments are also shared. Pre-training on code, before moving to mathematical data, enhances a model’s ability to address mathematical problems in scenarios both with and without access to tools. This provides new insight into the longstanding query: \textit{does code pre-training boost reasoning performance?} Our evidence suggests it does, at least within the mathematical domain.

    \item In contrast to common practice, we find that incorporating arXiv papers into the training corpus does not yield measurable benefits across the mathematical benchmarks considered in this study.
\end{itemize}

\textbf{Exploration and Analysis of Reinforcement Learning}

\begin{itemize}[topsep=0pt]
    \item We propose Group Relative Policy Optimization (GRPO), a reinforcement learning algorithm designed for both efficiency and effectiveness. Unlike traditional approaches such as Proximal Policy Optimization (PPO), GRPO eliminates the need for a critic by utilizing group score-based baseline estimation, which dramatically lowers the computational requirements during training.
    \item Our experiments reveal that applying GRPO to our instruction-tuned model, DeepSeekMath-Instruct, using only instruction-tuning data, yields substantial improvements in performance. Additionally, we find that this approach also benefits the model's out-of-domain capabilities throughout the reinforcement learning phase.
    \item We establish a cohesive framework for interpreting a range of techniques, including RFT, DPO, PPO, and GRPO. Our study includes systematic evaluations—comparing, for example, online to offline training, supervision of outcomes versus processes, and single-step against iterative reinforcement learning—to elucidate the critical features that define this unified paradigm.
    \item Leveraging insights from this unifying framework, we investigate why reinforcement learning is effective in these contexts and identify promising future research avenues to further enhance the reinforcement learning of large language models (LLMs).
\end{itemize}

\subsection{Summary of Evaluations and Metrics}

\begin{itemize}[topsep=0pt]
    \item \textbf{English and Chinese Mathematical Reasoning}: We perform thorough evaluations of our models across diverse English and Chinese datasets, spanning problem difficulty from elementary to collegiate levels. English benchmarks incorporate GSM8K \citep{gsm8k}, MATH \citep{MATH}, SAT \citep{llemma}, OCW Courses \citep{minerva}, and MMLU-STEM \citep{mmlu}. For Chinese, the benchmarks consist of MGSM-zh \citep{mgsm}, CMATH \citep{wei2023cmath}, Gaokao-MathCloze \citep{agieval}, and Gaokao-MathQA \citep{agieval}. We measure not only the capacity of the models to produce self-contained textual solutions without auxiliary tools but also their problem-solving proficiency using Python code.

    On the suite of English benchmarks, DeepSeekMath-Base shows comparable results to the proprietary Minerva 540B \citep{minerva}, and consistently outperforms all open-source foundation models—including Mistral 7B \citep{mistral} and Llemma-34B \citep{llemma}—regardless of math-specific pre-training, often by a large margin. Remarkably, DeepSeekMath-Base attains even higher accuracy on Chinese benchmarks, which may be attributed to our inclusion of high-quality multilingual pre-training data instead of restricting the dataset to English, as in prior studies \citep{minerva,llemma}. Through mathematical instruction tuning and reinforcement learning, both DeepSeekMath-Instruct and DeepSeekMath-RL exhibit exceptional outcomes, with DeepSeekMath-RL achieving over 50\% accuracy on the MATH dataset—representing the first such result for an open-source model in this domain.

\item \textbf{Formal Mathematics}: DeepSeekMath-Base is assessed on its capability for informal-to-formal theorem proving, following the setup in \citep{dsp_proof}, utilizing the miniF2F \citep{minif2f} dataset and the Isabelle \citep{isabelle} proof assistant. The model displays robust few-shot autoformalization performance.
\item \textbf{Natural Language Understanding, Reasoning, and Code}: To obtain a thorough evaluation of general language comprehension, reasoning, and programming abilities, we benchmark DeepSeekMath-Base against the Massive Multitask Language Understanding (MMLU) dataset \citep{mmlu}, comprising 57 multiple-choice topics spanning various fields, as well as BIG-Bench Hard (BBH) \citep{bbh}, a suite of 23 tasks focused predominantly on complex, multi-step reasoning, along with HumanEval \citep{codex} and MBPP \citep{mbpp}—both recognized as standard code model benchmarks. Mathematical pre-training is shown to enhance both the language understanding and reasoning proficiency of the model.
\end{itemize}

\section{Math Pre-Training}

\subsection{Data Collection and Decontamination} This section details the methodology used to construct the DeepSeekMath Corpus based on Common Crawl data. As illustrated in Figure \ref{fig:data_collect}, we employ an iterative data pipeline to systematically build an extensive mathematical text corpus from Common Crawl, starting with an initial seed dataset comprised of curated, high-quality math resources. Notably, the described pipeline is adaptable for other specialized domains, such as programming.

Initially, OpenWebMath \citep{openwebmath}—a dataset containing top-tier mathematical texts from the web—serves as our foundational seed corpus. A fastText model \citep{joulin2016fasttext} is then trained using this seed corpus to identify and retrieve additional mathematical web content with similar characteristics. For model training, we select 500,000 positive instances from the seed dataset and 500,000 negative samples from Common Crawl. An open-source implementation\footnote{\url{https://fasttext.cc}} is utilized with the following parameters: vector size of 256, learning rate of 0.1, a maximum n-gram length of 3, a minimum occurrence threshold of 3, and 3 epochs. To decrease the size of Common Crawl, we perform both URL deduplication and near-duplicate filtering, resulting in a collection of 40B HTML web pages. The trained fastText model is then employed to extract mathematically relevant pages from this deduplicated corpus. We score the retrieved pages using the fastText model’s outputs and retain only those with the highest scores, filtering out less relevant content. The amount of data retained is subsequently validated through pre-training runs on subsets comprising the top 40B, 80B, 120B, and 160B tokens, with the initial phase selecting the 40B top-scoring tokens for further processing.

Following the initial round of data gathering, a substantial portion of mathematical web pages remains uncollected, primarily due to the limited diversity within the fastText model's positive training set. To remedy this, we identify further mathematical domains to expand the seed corpus, thereby enhancing the effectiveness of the fastText classifier. Concretely, we first partition the entire Common Crawl dataset into non-overlapping domains, each defined by sharing an identical base URL. For every domain, we compute the fraction of web pages included during the first collection phase. If over 10\% of the domain’s pages are present in the initial data set, we classify it as math-related (e.g., \url{mathoverflow.net}). Next, we perform manual annotation of the URLs within these domains that correspond to mathematical content (e.g., \url{mathoverflow.net/questions}). Any remaining web pages linked to these annotated URLs, which were not previously collected, are incorporated into the seed set. This iterative approach broadens the range of positive samples, thus enabling the subsequent fastText models to better retrieve mathematical web content. After conducting four data collection iterations, we accumulate 35.5M mathematical web pages, comprising 120B tokens. During the fourth iteration, it becomes evident that nearly 98\% of the web pages were already obtained by the third iteration, prompting the termination of further collection efforts.

To minimize benchmark data leakage, we adopt the methodology from \cite{deepseek-coder} to eliminate web pages containing questions or solutions from English-language mathematics benchmarks such as GSM8K \citep{gsm8k} and MATH \citep{MATH}, as well as Chinese benchmarks like CMATH \citep{wei2023cmath} and AGIEval \citep{agieval}. Our filtering procedure removes any text segment that includes an exact match of a 10-gram string with any substring found within the benchmark datasets from our math corpus. In cases where the benchmark text is shorter than 10 grams but contains at least 3 grams, we rely on direct string matching to exclude the relevant contaminated pages.

\subsection{Validating the Quality of the DeepSeekMath~Corpus}

We conduct pre-training evaluations to compare the DeepSeekMath~Corpus against recently published mathematical text datasets:

\begin{itemize}[topsep=0pt]
    \item \textbf{MathPile} \citep{mathpile}: An 8.9B-token collection compiled from various sources, including textbooks, Wikipedia, ProofWiki, CommonCrawl, StackExchange, and arXiv—the latter making up more than 85\% of the dataset;
    \item \textbf{OpenWebMath} \citep{openwebmath}: A 13.6B-token dataset derived from filtering CommonCrawl for mathematical material;
    \item \textbf{Proof-Pile-2} \citep{llemma}: Composed of OpenWebMath, AlgebraicStack (10.3B tokens of mathematical code), and arXiv articles (28.0B tokens); in our experiments on Proof-Pile-2, we adhere to \cite{llemma}'s recommended source mixture ratio of 2:4:1 for arXiv, web, and code content respectively.
\end{itemize}

\subsubsection{Training Setting}
\label{sec:quality-policy}
We conduct specialized mathematical training on a foundational language model consisting of 1.3B parameters, structurally equivalent to the DeepSeek LLMs \citep{deepseek-llm}, referred to herein as DeepSeek-LLM 1.3B. For each mathematical dataset, we initiate an independent training process spanning 150B tokens. All experimental procedures utilize the resource-efficient HAI-LLM \citep{haillm} training platform. Adhering to the methodologies established for DeepSeek LLMs, optimization is carried out with the AdamW optimizer \citep{adamW}, using $\beta_1=0.9$, $\beta_2=0.95$, and $\mathrm{weight\_decay}=0.1$. The learning rate strategy consists of a multi-stage schedule: the learning rate ascends to its maximum over 2,000 warmup steps, diminishes to 31.6\% of this maximum after 80\% completion of training, and further drops to 10.0\% post 90\% of training. The maximum learning rate is maintained at 5.3e-4, with a token batch size of 4M and a context length capped at 4K tokens.

\subsubsection{Evaluation Results}

\textbf{The DeepSeekMath~Corpus demonstrates superior quality, multilingual breadth in mathematical content, and surpasses others in scale.}

\begin{itemize}[topsep=0pt]
    \item \textbf{High-quality}:
    Downstream effectiveness is measured across 8 mathematical evaluation datasets utilizing few-shot chain-of-thought prompting \cite{cot}. According to Table \ref{tab:corpora_comparison}, the model trained on the DeepSeekMath~Corpus consistently outperforms other models. Furthermore, Figure \ref{fig:corpora_comparisons} illustrates that after 50B tokens of training (equivalent to one complete epoch of Proof-Pile-2), DeepSeekMath-trained models outperform those trained on Proof-Pile-2, signaling that DeepSeekMath’s average data quality is notably higher.
    \item \textbf{Multilingual}:
    The DeepSeekMath~Corpus comprises multilingual content, with English and Chinese forming the bulk of the dataset. Results in Table \ref{tab:corpora_comparison} demonstrate that leveraging the DeepSeekMath~Corpus boosts mathematical reasoning for both English and Chinese prompts. In comparison, extant mathematical corpora focus mainly on English and often offer minimal gains—or even negative impact—when it comes to Chinese mathematical reasoning performance.
    \item \textbf{Large-scale}:
    Relative to previous mathematical datasets, the DeepSeekMath~Corpus is several-fold larger. As shown in Figure \ref{fig:corpora_comparisons}, training DeepSeek-LLM 1.3B on DeepSeekMath yields more rapid performance gains and sustained model improvement. On the other hand, baseline datasets are significantly smaller, require repeated epochs within the same training duration, and the resultant models plateau in capability much sooner.
\end{itemize}

\subsection{Training and Evaluating DeepSeekMath-Base 7B}

This section presents DeepSeekMath-Base 7B, a foundational model characterized by robust reasoning skills, particularly within mathematical domains. The model originates from DeepSeek-Coder-Base-v1.5 7B \citep{deepseek-coder} and undergoes training on 500B tokens. The composition of the training data is as follows: DeepSeekMath Corpus accounts for 56\%, AlgebraicStack constitutes 4\%, arXiv represents 10\%, 20\% is sourced from Github code, while the last 10\% comprises natural language text from Common Crawl in both English and Chinese. The majority of training procedures adhere to those outlined in Section \ref{sec:quality-policy}, though we specifically set the learning rate peak at 4.2e-4 and employ a batch size of 10 million tokens.

A thorough evaluation is conducted to investigate the mathematical problem-solving proficiency of DeepSeekMath-Base 7B. This evaluation examines its capability to independently generate comprehensive mathematical solutions, its facility with tool-assisted problem solving, and its performance in formal theorem proving. In addition, we report on the broader abilities of the model, including its proficiency in natural language understanding, logical reasoning, and programming.

\paragraph{Mathematical Problem Solving with Step-by-Step Reasoning}

DeepSeekMath-Base is assessed on its ability to address mathematical questions through few-shot chain-of-thought prompting \citep{cot} across eight datasets in both English and Chinese. These datasets include quantitative reasoning tasks (such as GSM8K \citep{gsm8k}, MATH \citep{MATH}, and CMATH \citep{wei2023cmath}) as well as multiple-choice evaluations (for example, MMLU-STEM \citep{mmlu} and Gaokao-MathQA \citep{agieval}), thereby spanning a broad spectrum of mathematical topics and problem complexities from elementary up to collegiate levels.

As illustrated in Table \ref{tab:base_math_cot}, DeepSeekMath-Base 7B consistently achieves the highest results among open-source base models across all eight evaluated benchmarks, outperforming models such as the widely adopted Mistral 7B \citep{mistral} and the recently introduced Llemma 34B \citep{llemma}, which was additionally trained on the Proof-Pile-2 dataset \citep{llemma} for mathematical tasks. Particularly, on the rigorous MATH competition dataset, DeepSeekMath-Base achieves a performance margin exceeding 10\% absolute improvement over existing open-source base models and surpasses Minerva 540B \citep{minerva}, a proprietary model 77 times larger that extends PaLM \citep{palm} and undergoes further mathematical pretraining.

\paragraph{Mathematical Problem Solving with Tool Use} To assess program-based mathematical reasoning, we employ few-shot program-of-thought prompting \citep{pot,pal} on both GSM8K and MATH datasets. In these experiments, models address each question by writing a Python program, utilizing packages like \textit{math} and \textit{sympy} for complex calculations, and the answer is taken as the output of the program execution. Table \ref{tab:base_math_pot_proof} demonstrates that DeepSeekMath-Base 7B outperforms the previous top-performing open-source model, Llemma 34B.

\paragraph{Formal Mathematics}

Formalization of mathematical proofs serves to verify their accuracy, improve reliability, and increase automation efficiency—areas receiving growing attention recently. We assess DeepSeekMath-Base 7B on informal-to-formal proof generation as described in \citep{dsp_proof}, requiring the translation of an informal statement and proof (with a corresponding formal version of the statement) into a formalized proof. Evaluation is conducted on miniF2F \citep{minif2f}, a formal Olympiad-level math benchmark, using few-shot prompts to generate proofs in Isabelle for each task. Following the protocol in \cite{dsp_proof}, models produce proof sketches, which are then completed with the Sledgehammer automated prover \citep{sledgehammer}. The results in Table \ref{tab:base_math_pot_proof} indicate that DeepSeekMath-Base 7B attains competitive performance in proof formalization.

\paragraph{Natural Language Understanding, Reasoning, and Code}

We examine the model’s performance on natural language understanding using MMLU \citep{mmlu}, on reasoning with BBH \citep{bbh}, and on programming via HumanEval \citep{codex} and MBPP \citep{mbpp}. According to Table \ref{tab:understanding_reasoning_code}, DeepSeekMath-Base 7B yields substantial improvements over its predecessor, DeepSeek-Coder-Base-v1.5 \citep{deepseek-coder}, for MMLU and BBH, signifying the beneficial effects of mathematical training on comprehension and reasoning tasks. Furthermore, through continued training that incorporates code tokens, DeepSeekMath-Base 7B preserves the coding abilities of DeepSeek-Coder-Base-v1.5 on HumanEval and MBPP. On the whole, DeepSeekMath-Base 7B also shows marked superiority over the general-purpose Mistral 7B \citep{mistral} across the evaluated reasoning and coding benchmarks.

\section{Supervised Fine-Tuning}

\subsection{SFT Data Curation}

We compile a mathematical instruction-tuning dataset encompassing both English and Chinese language questions from a broad spectrum of mathematical domains and difficulty levels. Each problem is paired with solutions formatted in chain-of-thought (CoT) \citep{cot}, program-of-thought (PoT) \citep{pot,pal}, or tool-integrated reasoning style \citep{tora}. The full dataset comprises 776,000 training instances.

\begin{itemize}[topsep=0pt]
    \item \textbf{English mathematical datasets}: We provide tool-integrated solution annotations for problems from GSM8K and MATH, and utilize a selected subset of MathInstruct \citep{MathInstruct} as well as the training portion of Lila-OOD \citep{lila}, featuring solutions in CoT or PoT formats. This collection spans a variety of mathematical disciplines such as algebra, probability, number theory, calculus, and geometry.
    \item \textbf{Chinese mathematical datasets}: Our collection includes Chinese K-12 mathematics problems distributed across 76 different sub-fields, such as linear equations, with annotated solutions following both CoT and tool-integrated formats.
\end{itemize}

\subsection{Training and Evaluating DeepSeekMath-Instruct 7B}

This section presents DeepSeekMath-Instruct 7B, which is further instruction-tuned for mathematical tasks using DeepSeekMath-Base as the foundation. During training, multiple examples are concatenated at random, constrained by a 4,000-token maximum context window. The model is optimized for 500 steps with a batch size set to 256 and a constant learning rate of $5 \times 10^{-5}$.

We assess mathematical proficiency of the models, both in settings without tool access and when tools are available, across four quantitative reasoning benchmarks in English and Chinese. Comparative evaluations are performed against top-performing contemporaneous models, specifically:

\begin{itemize}[topsep=0pt]
    \item \textbf{Closed-source models} include: (1) the GPT suite, notably GPT-4 \citep{gpt4} and GPT-4 Code Interpreter~\footnote{\url{https://openai.com/blog/chatgpt-plugins\#\#code-interpreter}}, (2) Gemini Ultra and Gemini Pro \citep{gemini}, (3) Inflection-2 \citep{inflection-2}, (4) Grok-1~\footnote{\url{https://x.ai/model-card}}, and (5) Baichuan-3~\footnote{\url{https://www.baichuan-ai.com}}, together with (6) the newest release from the GLM family, GLM-4~\footnote{\url{https://open.bigmodel.cn/dev/api\#glm-4}} \citep{glm}.
\end{itemize}

These models serve general purposes and have typically been subject to various alignment strategies.
    \item \textbf{Open-source models} comprise: general-purpose models such as (1) DeepSeek-LLM-Chat 67B \citep{deepseek-llm}, (2) Qwen 72B \citep{qwen}, (3) SeaLLM-v2 7B \citep{seallm}, and (4) ChatGLM3 6B \citep{chatglm3}; in addition to models specifically enhanced for mathematics: (5) InternLM2-Math 20B~\footnote{\url{https://github.com/InternLM/InternLM-Math}}, which extends InternLM2 by including mathematical pretraining and subsequent instruction tuning, (6) Math-Shepherd-Mistral 7B, which leverages PPO optimization \citep{schulman2017proximal} on Mistral 7B \citep{mistral} utilizing a reward model guided by process supervision, (7) the WizardMath family \citep{wizardmath}, which bolsters mathematical reasoning capabilities in both Mistral 7B and Llama-2 70B \citep{llama2} via evolve-instruct (a variant of instruction tuning employing AI-generated instructions) and PPO optimization on datasets dominated by GSM8K and MATH, (8) MetaMath 70B \citep{metamath}, where Llama-2 70B is adapted with an expanded GSM8K and MATH corpus, (9) ToRA 34B \cite{tora}, a CodeLlama 34B model tailored for tool-augmented mathematical reasoning, and (10) MAmmoTH 70B \citep{MathInstruct}, consisting of Llama-2 70B refined via instruction tuning using MathInstruct.

Referencing Table \ref{tab:sft_rl_math}, within the evaluation framework prohibiting the use of tools, DeepSeekMath-Instruct 7B exhibits outstanding capability in stepwise reasoning. Remarkably, on the advanced MATH benchmark, this model achieves at least a 9\% absolute margin over all other open-source counterparts and most commercial systems (including Inflection-2 and Gemini Pro). This holds even when compared with considerably larger architectures (such as Qwen 72B) or those specifically refined with math-oriented RL methodologies (e.g., WizardMath-v1.1 7B). DeepSeekMath-Instruct’s performance matches that of leading Chinese proprietary systems like GLM-4 and Baichuan-3 on MATH, but it still does not reach the proficiency levels of GPT-4 and Gemini Ultra.

In the scenario where models are permitted to combine natural language reasoning with code-based tool usage to address problems, DeepSeekMath-Instruct 7B attains close to 60\% accuracy on MATH, thereby outperforming every publicly available open-source system. Across other benchmarks, its results are on par with DeepSeek-LLM-Chat 67B, the former state-of-the-art, despite the latter's tenfold increase in size.

\section{Reinforcement Learning}

\subsection{Group Relative Policy Optimization} Reinforcement learning (RL) has demonstrated substantial utility for further boosting LLMs’ mathematical reasoning after completing the Supervised Fine-Tuning (SFT) phase \citep{wang2023math,wizardmath}. We present here our approach—Group Relative Policy Optimization (GRPO)—an RL algorithm designed to be both resource-efficient and effective.

\subsubsection{From PPO to GRPO} Proximal Policy Optimization (PPO) \citep{schulman2017proximal} stands as a prominent actor-critic RL method, extensively utilized during the RL-based fine-tuning of LLMs \citep{ouyang2022training}. PPO seeks to enhance model behavior by maximizing the following objective:

\begin{equation}
\footnotesize
    \mathcal{J}_{PPO}(\theta) = \mathbb{E}{[q \sim P(Q), o \sim \pi_{\theta_{old}}(O|q)]} \frac{1}{|o|} \sum_{t=1}^{|o|} \min \left[ \frac{\pi_\theta(o_{t} | q, o_{<t})}{\pi_{\theta_{old}}(o_{t} | q, o_{<t})} A_{t}, \text{clip} \left( \frac{\pi_\theta(o_{t} | q, o_{<t})}{\pi_{\theta_{old}}(o_{t} | q, o_{<t})}, 1 - \varepsilon, 1 + \

Here, $\pi_{\theta}$ and $\pi_{\theta_{old}}$ denote the policy models currently being trained and the previous iteration, respectively. The variables $q$ and $o$ refer to questions and generated outputs, which are sampled from the question dataset and the old policy $\pi_{\theta_{old}}$. The hyper-parameter $\varepsilon$ plays a crucial role in PPO by providing a clipping mechanism to enhance training stability. The term $A_t$ indicates the advantage, which is estimated by employing Generalized Advantage Estimation (GAE) \citep{gae}, using the reward sequence $\{r_{\ge t}\}$ along with a value function $V_{\psi}$ learned during training. Consequently, PPO necessitates that a value function be learned in parallel with the policy network. To mitigate the risk of excessively optimizing the reward model, a typical method is to introduce a per-token KL-divergence penalty—relative to a reference model—directly into the reward for each token, as in \citep{ouyang2022training}:

\begin{equation}
    r_{t} = r_\varphi(q, o_{\le t}) - \beta \log\frac{\pi_{\theta}(o_{t}|q, o_{<t})}{\pi_{ref}(o_{t}|q, o_{<t})},
\label{eq:PPO-reward}
\end{equation}

where $r_\varphi$ is the learned reward model, $\pi_{ref}$ stands for the reference policy (commonly the initial SFT model), and $\beta$ denotes the weighting of the KL regularization.

Since the value function utilized in PPO generally matches the policy model in size, it incurs significant memory and computational overhead. Further, when training with reinforcement learning, the value function is applied as a baseline in advantage calculations to control variance. In LLM-based RL settings, the reward model often provides feedback solely for the final token, complicating the construction of a reliable, token-level value function. To resolve this, we introduce Group Relative Policy Optimization (GRPO) as depicted in Figure \ref{fig:grpo}. GRPO eliminates the need for training an additional value function, as is required in PPO, and substitutes it with the group-average reward from multiple sampled responses to the same prompt, using this mean as the baseline. Explicitly, for a given question $q$, GRPO draws a group of outputs $\{o_1, o_2, \cdots, o_G\}$ from the old policy $\pi_{\theta_{old}}$, and trains the policy by optimizing the objective below:

\begin{equation}
\footnotesize
\begin{split}
    \mathcal{J}_{GRPO}(\theta) &= \mathbb{E}{[q \sim P(Q), \{o_i\}_{i=1}^G \sim \pi_{\theta_{old}}(O|q)]}  \\
    & \frac{1}{G}\sum_{i=1}^G\frac{1}{|o_i|} \sum_{t=1}^{|o_i|} \left\{ \min \left[ \frac{\pi_\theta(o_{i,t} | q, o_{i,<t})}{\pi_{\theta_{old}}(o_{i,t} | q, o_{i,<t})} \hat{A}_{i,t}, \text{clip} \left( \frac{\pi_\theta(o_{i,t} | q, o_{i,<t})}{\pi_{\theta_{old}}(o_{i,t} | q, o_{i,<t})}, 1 - \varepsilon, 1 + \varepsilon \right)  \hat{A}_{i,t} \right] - \beta \mathbb{D}_{KL}\left[\pi_{\theta} || \pi_{ref}\right]\right\} ,
\end{split}
\label{eq:GRPO-obj}
\end{equation}

Within this formulation, $\varepsilon$ and $\beta$ remain as hyper-parameters, while $\hat{A}_{i,t}$ represents the advantage term, computed only from the relative rewards within the output group—a method further explained in later sections. By using group-based relative calculations for advantage, GRPO naturally fits with the comparative evaluation approach of reward models, which are commonly trained on pairwise or listwise assessments of answers to the same query. Importantly, whereas the KL penalty in (\ref{eq:PPO-reward}) modifies the reward function itself, GRPO integrates the KL divergence between the updated policy and the reference policy directly into the loss, thus simplifying the calculation of $\hat{A}_{i,t}$. In addition, rather than employing the

\begin{algorithm}[t]
  \small
  \caption{Iterative Group Relative Policy Optimization}
  \textbf{Input} initial policy model $\pi_{\theta_{\text{init}}}$; reward models $r_\varphi$; task prompts $\mathcal{D}$; 
  hyperparameters $\varepsilon$, $\beta$, $\mu$
  \begin{algorithmic}[1]
    \State Set policy model $\pi_\theta \leftarrow \pi_{\theta_{\text{init}}}$
    \For{iteration = 1, \dots, I}
       \State Assign reference model $\pi_{ref} \leftarrow \pi_{\theta}$
      \For{step = 1, \dots, M}
      \State Draw a minibatch $\mathcal{D}_b$ from $\mathcal{D}$
      \State Save old policy $\pi_{\theta_{old}} \leftarrow \pi_{\theta}$ 
      \State For each prompt $q \in \mathcal{D}_b$, sample $G$ completions $\{o_i\}_{i=1}^G \sim \pi_{\theta_{old}} (\cdot \mid q) $
      \State Evaluate rewards $\{r_i\}_{i=1}^{G}$ for all outputs $o_i$ using $r_{\varphi}$
      \State For each $t$-th token of $o_i$, estimate $\hat{A}_{i,t}$ via group-relative advantage computation
      \For{GRPO iteration = 1, \dots, $\mu$}
        \State Update $\pi_{\theta}$ by maximizing the group-relative policy objective (Equation \ref{eq:GC-GRPO})
      \EndFor
    \EndFor 
    \State Continue training $r_\varphi$ using a replay buffer. 
    \EndFor 
  \end{algorithmic}
  \textbf{Output} $\pi_\theta$
  \label{alg:iter-grpo}
\end{algorithm}

\subsubsection{Outcome Supervision RL with GRPO} In this framework, for each input question $q$, a set of candidate responses $\{o_1, o_2, \cdots, o_G\}$ is sampled from the previous policy $\pi_{\theta_{old}}$. Each output is assigned a score via a reward model, resulting in the reward vector $\mathbf{r}=\{r_1, r_2, \cdots, r_G\}$. To enable effective credit assignment, these scores are standardized by centering and scaling using the group's mean and standard deviation. Under outcome supervision, only the normalized terminal reward for each completion $o_i$ is supplied, and this value is allocated as the advantage $\hat{A}_{i, t}$ to all tokens within the output: $\hat{A}_{i, t} = \widetilde{r}_i = \frac{r_i- {\rm mean}(\mathbf{r})}{{\rm std}(\mathbf{r})}$. The policy is subsequently refined by maximizing the objective given in equation (\ref{eq:GRPO-obj}).

\subsubsection{Process Supervision RL with GRPO} Relying solely on outcome rewards, which provide feedback at the sequence level, may fail to sufficiently guide the policy, especially in complex reasoning tasks. In alignment with \cite{wang2023math}, we employ process supervision, which delivers feedback after each reasoning stage. More concretely, given an input $q$ and generated samples $\{o_1, o_2, \cdots, o_G\}$, a process-oriented reward model supplies step-level rewards: $\mathbf{R} = \{ \{r_1^{index(1)},\cdots,r_1^{index(K_1)}\}, \cdots,  \{r_G^{index(1)},\cdots,r_G^{index(K_G)}\} \}$, where $index(j)$ indicates the terminal token of the $j$-th reasoning step and $K_i$ counts the total reasoning steps for output $o_i$. These stepwise rewards are normalized according to the overall mean and standard deviation: $\widetilde{r}_i^{index(j)} = \frac{r_i^{index(j)} - {\rm mean(\mathbf{R})}}{{\rm std(\mathbf{R})}}$.
The advantage assigned to each token is then computed as the cumulative sum of normalized rewards for all steps ending at or after position $t$, i.e., $\hat{A}_{i, t} = \sum_{index(j) \ge t} \widetilde{r}_i^{index(j)}$.
Finally, the policy is updated by maximizing

\subsubsection{Iterative RL with GRPO} During reinforcement learning, the previously trained reward model may become inadequate for effectively guiding the evolving policy model. To address this, we experiment with an iterative approach to RL using GRPO. As described in Algorithm \ref{alg:iter-grpo}, the iterative GRPO procedure involves generating updated datasets for the reward model through samples drawn from the policy model. The reward model is then continually refined utilizing a replay buffer that includes 10\% of past data. Subsequently, the updated policy model is assigned as the new reference model, and the policy is further optimized using the latest version of the reward model.

\subsection{Training and Evaluating DeepSeekMath-RL}

For reinforcement learning, we employ the DeepSeekMath-Instruct 7B as the base model. The RL training corpus comprises approximately 144K chain-of-thought questions on GSM8K and MATH, sourced from the SFT dataset. To specifically analyze the effect of RL on benchmarks that are unseen during RL, other SFT examples are omitted. Construction of the reward model's training set adheres to the methodology in \citep{wang2023math}. The initial reward model is trained on DeepSeekMath-Base 7B with a learning rate of 2e-5. The GRPO policy model employs a learning rate of 1e-6 and a KL coefficient of 0.04. Each input question is used to produce $64$ samples, with a maximum sequence length of 1024 and a batch size of 1024. The policy model undergoes a single parameter update after each exploration round. For evaluation, DeepSeekMath-RL 7B is benchmarked using the same benchmarks as DeepSeekMath-Instruct 7B. In the context of DeepSeekMath-RL 7B, tasks involving GSM8K and MATH with chain-of-thought prompts are considered in-domain, while all remaining benchmarks are categorized as out-of-domain.

Table \ref{tab:sft_rl_math} presents a comparison of open-source and closed-source models evaluated with both chain-of-thought and tool-integrated reasoning strategies on English and Chinese tasks. The main observations are as follows: (1) DeepSeekMath-RL 7B achieves accuracies of 88.2\% on GSM8K and 51.7\% on MATH using chain-of-thought approaches. This performance exceeds all open-source competitors within the 7B--70B parameter range and surpasses most closed-source models as well. (2) Notably, DeepSeekMath-RL 7B’s training is limited to chain-of-thought-style instruction tuning data drawn only from GSM8K and MATH, using DeepSeekMath-Instruct 7B as its base model. Despite this narrow training data, DeepSeekMath-RL 7B consistently outperforms DeepSeekMath-Instruct 7B in every evaluation category, highlighting the strong impact of reinforcement learning.

\section{Discussion}
Here, we summarize insights gained from our pre-training and reinforcement learning experiments.

\subsection{Lessons Learnt in Pre-Training}

We begin by sharing insights regarding pre-training. Unless specified otherwise, the training protocols adhere to the descriptions in Section \ref{sec:quality-policy}. The term "DeepSeekMath Corpus" in this discussion refers to an 89-billion-token collection assembled in the second data iteration.

\subsubsection{Code Training Benefits Mathematical Reasoning}

Although it has often been suggested but not systematically confirmed that training on code enhances reasoning abilities, our work attempts to clarify this claim—at least with regard to mathematical problem solving. Our findings indicate that code training bolsters a model’s mathematical reasoning skills, with or without the use of external tools.

To investigate the relationship between code training and mathematical reasoning, we adopted both two-stage and one-stage training schemes, described as follows:

\textbf{Two-Stage Training}

\begin{itemize}[topsep=0pt]
    \item \textbf{Code Training for 400B Tokens $\rightarrow$ Math Training for 150B Tokens}: DeepSeek-LLM 1.3B is initially trained on 400B tokens sourced from code, followed by an additional 150B tokens from mathematics datasets;
    \item \textbf{General Training for 400B Tokens $\rightarrow$ Math Training for 150B Tokens}: For comparison, we perform another training regimen where general-domain tokens (sampled from DeepSeek-AI’s comprehensive general-purpose corpus) replace code tokens in the first stage, aiming to discern whether code-specific pretraining yields superior enhancements to mathematical reasoning compared to broad-domain data.
\end{itemize}

\textbf{One-Stage Training}

\begin{itemize}[topsep=0pt]
    \item \textbf{Math Training for 150B Tokens}: The model DeepSeek-LLM 1.3B undergoes a single phase of training solely on 150B tokens of mathematical content;
    \item \textbf{Training on a mixture of 400B Code Tokens and 150B Math Tokens}: Because sequential math training after code training diminishes code capability, we further analyze the impact of a joint training regime—using a blended dataset of 400B code tokens and 150B math tokens—on both mathematical reasoning and the retention of code proficiency, seeking to avoid catastrophic forgetting.
\end{itemize}

\paragraph{Results}
Table \ref{tab:code-math} and Table \ref{tab:code-math-general-eval} present evaluation outcomes across varying training protocols.

Pretraining on code data notably enhances performance on program-assisted mathematical tasks under both staged and joint training frameworks. Specifically, Table \ref{tab:code-math} shows that models trained first on code exhibit marked improvements in solving GSM8K and MATH benchmarks with Python, even before subsequent mathematical specialization; additional math-specific training provides further gains. When mixing code and math tokens in a unified training phase, this strategy not only helps counteract catastrophic forgetting typically observed with multi-phase setups but also augments both code performance (see Table \ref{tab:code-math-general-eval}) and the ability to leverage program-aided reasoning (refer to Table \ref{tab:code-math}).

Further, code-centric pretraining also offers advantages for mathematical reasoning in settings that do not rely on external tools. Within the two-phase scheme, initial code training confers moderate gains and facilitates more effective subsequent math specialization, collectively producing optimal results. In contrast, joint exposure to code and math data in a single phase appears detrimental to tool-free mathematical reasoning. A plausible hypothesis for this phenomenon is that DeepSeek-LLM 1.3B’s restricted parameter budget hampers its ability to simultaneously internalize both domains to full effect.

\subsubsection{Limited Effectiveness of ArXiv Papers in Enhancing Mathematical Reasoning}

Despite their widespread use in pre-training datasets for mathematical language models \citep{minerva,gpt-f,llemma,mathpile}, the specific contribution of arXiv papers to mathematical reasoning ability has not been comprehensively assessed. Surprisingly, our experimental findings indicate that incorporating arXiv papers yields negligible improvement in models’ mathematical reasoning capabilities. We investigated this effect across different model architectures, including DeepSeek-LLM 1.3B and DeepSeek-Coder-Base-v1.5 7B \citep{deepseek-coder}, utilizing distinct arXiv-based datasets processed through varying methodologies:

\begin{itemize}[topsep=0pt]
    \item \textbf{MathPile} \citep{mathpile}: An 8.9B-token dataset curated with heuristic rules for cleaning and filtering, consisting of over 85\% scientific arXiv papers.
    \item \textbf{ArXiv-RedPajama} \citep{redpajama}: A 28.0B-token corpus of arXiv LaTeX content, stripped of preambles, comments, macros, and bibliographies.
\end{itemize}

For evaluation, we trained DeepSeek-LLM 1.3B on 150B tokens and DeepSeek-Coder-Base-v1.5 7B on 40B tokens using each arXiv dataset. Our results consistently show that exposure to arXiv content does not translate into notable gains in mathematical reasoning performance. In certain cases, models exhibited stagnation or even declines in accuracy on several mathematics benchmarks of varied difficulty covered in this work, including datasets focused on quantitative problem solving such as GSM8K and MATH (refer to Table \ref{tab:arxiv-cot}), multiple-choice evaluations like MMLU-STEM (Table \ref{tab:arxiv-cot}), and formal mathematics assessment via miniF2F (Table \ref{tab:arxiv_atp}).

It is important to highlight the preliminary nature of this conclusion, due to several unaddressed factors:

\begin{itemize}[topsep=0pt]
    \item Potential impact of arXiv-derived data on specialized mathematical subtasks not represented here, such as the process of transforming formal theorem statements or proofs into their informal equivalents;
    \item Effects that may arise when arXiv content is integrated alongside other data modalities;
    \item The possibility that beneficial effects could surface with increased model scale.
\end{itemize}

As a result, additional investigation is warranted to address these open questions, which we defer to subsequent research.

\subsection{Perspectives on Reinforcement Learning}
\subsubsection{Unifying Training Methodologies} This section introduces a unified analytical framework for a variety of training algorithms—SFT, RFT, DPO, PPO, GRPO—and presents empirical results to dissect the underlying components of this unified perspective. Typically, for a given method, the parameter gradient with respect to $\theta$ can be formulated as follows:

\begin{equation}
    \nabla_{\theta}\mathcal{J}_{\textcolor{red}{\mathcal{A}}}(\theta) = \mathbb{E}[\underbrace{(q,o) \sim \textcolor{red}{\mathcal{D}}}_{Data \ Source}]\left( \frac{1}{|o|} \sum_{t=1}^{|o|}  \underbrace{GC_{{\mathcal{A}}}(q, o, t, \textcolor{red}{\pi_{{rf}}})}_{Gradient \ Coefficient}  \nabla_{\theta}\log \pi_{\theta}(o_t | q, o_{<t})\right).
\label{eq:objective}
\end{equation}

This formulation comprises three fundamental elements: 1) \textit{Data Source $\mathcal{D}$}, which specifies the distribution of samples utilized for model training; 2) \textit{Reward Function $\pi_{{rf}}$}, serving as the origin of the feedback signal for learning; and 3) \textit{Algorithm $\mathcal{A}$}, which computes the gradient coefficient $GC$ by integrating the training data and reward signals, thereby adjusting the model parameters via reinforcement or penalization. Within this unified framework, we examine a range of prominent approaches:

\begin{itemize}
    \item  \textbf{Supervised Fine-tuning (SFT)}: This method involves further training a pretrained model on curated datasets composed of human-annotated SFT examples.
    \item \textbf{Rejection Sampling Fine-tuning (RFT)}: In RFT, the SFT model is additionally fine-tuned using filtered response samples produced by the SFT model for SFT-related questions, where responses are selected based on their correctness.
    \item \textbf{Direct Preference Optimization (DPO)}: DPO enhances the SFT model by applying fine-tuning with an augmented dataset consisting of paired outputs from the SFT model, optimized via a pairwise preference-based DPO loss.
    \item \textbf{Online Rejection Sampling Fine-tuning (Online RFT)}: Unlike standard RFT, Online RFT commences from the SFT model but incrementally updates the policy by sampling and training on new outputs from the evolving real-time policy model.
    \item \textbf{PPO/GRPO}: With PPO/GRPO, the policy model begins as the SFT model and is further strengthened by sampling outputs from the ongoing policy and applying reinforcement learning procedures.
\end{itemize}

The individual elements of these approaches are outlined in Table \ref{tab:data_policy}. For a more thorough exposition of the derivation, readers are directed to Appendix \ref{app:analysis-rl}.

\paragraph{Observation about Data Source} We categorize data sources as either online or offline sampling. In the context of online sampling, training data is obtained directly from the ongoing interactions produced by the policy model during training, whereas offline sampling leverages data drawn from the pre-existing SFT model's outputs. Notably, both RFT and DPO operate within the offline regime, while Online RFT and GRPO utilize data sampled through online means.

Referencing Figure \ref{fig:rl-analysis}, our analysis reveals that Online RFT demonstrates a marked improvement over RFT across two evaluated benchmarks. In the initial phase of training, Online RFT and RFT exhibit nearly equivalent performance. However, as training progresses, Online RFT achieves a clear and growing performance edge, showcasing the benefits of employing an online training approach. This observation aligns with intuition: at the outset, the policy and the SFT model are nearly identical, yielding only slight differences in sampled outputs. As training advances, the policy diverges further, and samples generated in real-time reflect more pronounced distinctions, with the online sampling mechanism conferring notable advantages.

\paragraph{Observation about Gradient Coefficient} The transformation of the reward signal into the gradient coefficient is fundamental for parameter updates within these algorithms. In our experimental setup, we distinguish between two reward functions: `Rule,' which evaluates the accuracy of a response, and `Model,' where a dedicated reward model, trained using rule-based labels, assigns a score to each output. As described by Equations \ref{eq:GC-RFT} and \ref{eq:GC-GRPO}, a principal distinction emerges: GRPO modifies its gradient coefficient in direct response to the reward magnitude from the reward model, enabling more granular encouragement or discouragement based on the relative merit of different responses. Conversely, Online RFT applies equal reinforcement to all correctly answered responses and provides no penalization for incorrect ones, thereby lacking the differential weighting found in GRPO.

As illustrated in Figure \ref{fig:rl-analysis}, GRPO outperforms online RFT, thereby underscoring the effectiveness of modifying positive and negative gradient coefficients. Moreover, the GRPO+PS configuration exhibits better results than GRPO+OS, which highlights the advantage of implementing step-aware, fine-grained gradient coefficients. We further investigate the impact of iterative RL by running two rounds of iteration in our experimental setup. According to Figure \ref{fig:iter-rl}, it is evident that applying iterative RL yields considerable gains, with the most pronounced improvement observed during the first iteration.

\subsubsection{Why RL Works?} This study applies reinforcement learning to a subset of instruction tuning data and observes substantial gains over the baseline instruction-tuned model. To clarify the mechanism behind RL’s efficacy, we assess both Pass@K and Maj@K metrics for the Instruct and RL-optimized models across two benchmarks. Figure \ref{fig:maj-pass} demonstrates that RL substantially boosts Maj@K performance while leaving Pass@K largely unaffected. This suggests that RL enhances robustness in the model’s output distribution, implying that the observed gains stem primarily from promoting correct responses within the TopK outputs, rather than advancing intrinsic model competence. Likewise, \citep{wang2023making} discuss a \textbf{misalignment problem} in reasoning tasks when using the SFT model, showing that reasoning capabilities can be improved through preference alignment approaches \citep{yuan2023rrhf,song2023preference,wang2023making}.

\subsubsection{How to Achieve More Effective RL?}
\label{sec:effec_rl}
Our results affirm that RL is particularly effective on mathematical reasoning problems. We additionally propose a unified framework for interpreting various prominent training approaches, situating them as variants or simplifications of RL methodologies. As formalized in Equation \ref{eq:objective}, the framework revolves around three main elements: Data Source, Algorithm, and Reward Function. We suggest prospective directions for improvement concerning each of these aspects.

\paragraph{Data Source} The data source constitutes the foundational material for any training approach. Specifically, within RL, this refers to the collection of unlabeled questions, paired with outputs sampled from the policy model. For our experiments, we limited ourselves to question data utilized during instruction tuning and adopted simple nucleus sampling for response generation. We surmise that this restricted data and sampling approach may account for why our RL process primarily enhances Maj@K outcomes. Going forward, we aim to extend our RL approach to encompass out-of-distribution prompts, supplemented by \textbf{advanced sampling (decoding) strategies}—such as those employing tree-search techniques \citep{yao2023tree}. In addition, the implementation of \textbf{efficient inference techniques} \citep{xia-etal-2023-speculative,leviathan2023fast,kwon2023efficient,xia2024unlocking}—which dictate how effectively policy models can explore their space—will be critically important.

\paragraph{Algorithms} Algorithms leverage both data and reward signals to compute gradients, which are then used to adjust model parameters. As indicated in Equation \ref{eq:objective}, contemporary approaches predominantly \textbf{TRUST} the feedback from the reward function when deciding to increase or decrease the conditional probability assigned to specific tokens. Nevertheless, the reliability of this reward signal cannot be guaranteed, particularly in highly intricate problem settings. For instance, despite rigorous annotation by expert annotators, the PRM800K datasets \citep{lightman2023let} still contain about 20\% incorrect labels\footnote{\url{https://github.com/openai/prm800k/issues/12\#issuecomment-1728491852}}. In response to these challenges, our investigation centers on reinforcement learning techniques capable of maintaining robustness in the presence of noisy reward feedback. We propose that alignment strategies under the \textbf{WEAK-TO-STRONG} framework \citep{burns2023weak} could represent a transformative step in the development of future learning algorithms.

\paragraph{Reward Function} The reward function acts as the principal mechanism for supplying training feedback. Within RL paradigms, this typically corresponds to a neural reward model. Our analysis identifies three pivotal avenues for advancing reward model research: 1) \textbf{Strategies for improving reward model generalization.} Ensuring robust generalization enables the reward model to evaluate out-of-distribution queries and sophisticated decoding results. Absent this property, reinforcement learning may simply stabilize the output space of LLMs without delivering substantive improvement; 2) \textbf{Capturing reward model uncertainty.} Effectively modeling uncertainty may facilitate integration between imprecise reward models and weak-to-strong training paradigms; 3) \textbf{Constructing high-quality process reward models efficiently}, which are able to deliver detailed and actionable signals throughout complex reasoning chains \citep{lightman2023let,wang2023math}.

\section{Conclusion, Limitation, and Future Work}

In this paper, we introduce DeepSeekMath, which surpasses all existing open-source models on the competition-level MATH benchmark and attains performance that is close to that of proprietary models. The foundation of DeepSeekMath~is DeepSeek-Coder-v1.5 7B, which is subsequently continually trained for 500B tokens; a substantial portion of this data comprises 120B math-specific tokens extracted from Common Crawl. Through a comprehensive ablation analysis, we identify that mathematical data from web sources holds considerable value, whereas data from arXiv appears less advantageous than anticipated. We also propose Group Relative Policy Optimization (GRPO), an adaptation of Proximal Policy Optimization (PPO), capable of significantly enhancing mathematical reasoning performance while reducing memory requirements. Experimental outcomes demonstrate that GRPO provides notable gains, even when DeepSeekMath-Instruct 7B already achieves strong results on standard benchmarks. Additionally, we outline a unified perspective to contextualize multiple techniques and highlight a set of prospective pathways for advancing reinforcement learning methods.

Despite DeepSeekMath's strong results on quantitative reasoning tasks, it exhibits comparatively weaker performance in areas such as geometry and theorem proving when measured against proprietary systems. For example, in preliminary assessments, the model fails to solve problems related to triangles and ellipses, likely due to biases in data selection during both pre-training and fine-tuning phases. Model size is another constraint—DeepSeekMath~underperforms GPT-4 in scenarios that leverage few-shot learning: while GPT-4 shows marked improvement with few-shot prompts, DeepSeekMath~yields similar results regardless of zero-shot or few-shot evaluation. Looking ahead, we aim to enhance our data curation strategies to build even higher-quality pre-training datasets. Furthermore, we plan to investigate additional approaches (see Section \ref{sec:effec_rl}) to develop more effective reinforcement learning methodologies for large language models.

\end{document}